\documentclass[10pt]{article}
\usepackage[margin=20mm]{geometry}
\usepackage{amsmath,amssymb,mathtools,fontspec,unicode-math,xurl,hyperref}
\IfFontExistsTF{DejaVu Serif}{\setmainfont{DejaVu Serif}}{\setmainfont{Latin Modern Roman}}
\setmathfont{latinmodern-math.otf}
\setlength{\parindent}{0pt}\setlength{\parskip}{6pt}\sloppy
\title{Initial Scientific Review of
Fast Randomized Least-Squares Solvers Can Be
Just as Accurate and Stable as Classical Direct Solvers}
\author{Bounded scientific review: M25-0427}
\date{}
\begin{document}
\maketitle
Abstract: The general-preconditioner extension, Theorem 6.10, is false with its stated scale-independent inner-error hypothesis. An explicit floating-point family satisfies all of its assumptions but fails both asserted stability conclusions. This Major finding is distinct from the normalized sketch-QR theory. A second Major finding concerns the recommended printed FOSSILS algorithm, whose reversed parameter ratio and other inconsistencies are absent from the inspected pinned code. Two Minor findings concern the condition-number constant in the Lanczos proof and its zero-right-hand-side guard. The positive normalized proof mechanism and the limits of the CG/LSQR claims are recorded. Scientific adjudication remains pending.\par
\section*{Identity and scope}
The source is E. N. Epperly, M. Meier and Y. Nakatsukasa, Fast randomized least-squares solvers can be just as accurate and stable as classical direct solvers, DOI https://doi.org/10.1002/cpa.70013, acquired as arXiv:2406.03468v3 [source]. The PDF carries the arXiv date 20 August 2025 and a manuscript date of 21 August 2025. Its 45 pages have SHA256 `ee2f3d97793fbc72cd32f4f09cccba96e39a4a9b8415bba108a1fcde2bebbfdd`. All source text and appendices were read. Numbered locators refer to this acquired version, with matching physical and printed pages; equivalence with a journal version is not asserted.\par
The bounded preserved-reference search recovered no substantive predecessor report for this paper. This is an initial review. The retained repository [code] is pinned to `afc7c8c071e4347c22b729c52c8492aa5678f009`. The files `code/fossils.m` and `code/metasolver.m` were read statically in full; no author code or numerical benchmark was run. Separate reviewer calculations and actual source-page inspections support the findings.\par
\section*{The normalized theory and its hypotheses}
The problem is full-column-rank least squares, \(x=\arg\min_y\left\lVert b-Ay\right\rVert_2\). The analysis normalizes \(\left\lVert A\right\rVert=\left\lVert b\right\rVert=1\), assumes \(\kappa u\ll1\), and suppresses dimension-polynomial constants in \(\lesssim\). For a sketch satisfying \[(1-\eta)\left\lVert Ay\right\rVert\le\left\lVert SAy\right\rVert\le(1+\eta)\left\lVert Ay\right\rVert,
 \qquad SA=QR,\] the singular values of \(AR^{-1}\) lie between \((1+\eta)^{-1}\) and \((1-\eta)^{-1}\). This controls absolute scale as well as condition number.\par
With \(M=\widehat R^{-T}A^TA\widehat R^{-1}\), the refinement forms \(c_i=\widehat R^{-T}A^T(b-A\widehat x_i)\), approximately solves \(M\delta y_i=c_i\), and updates \(\widehat x_{i+1}=\widehat x_i+\widehat R^{-1}\delta y_i\). Theorem 6.1 requires the absolute inner-error estimate \[\left\lVert\operatorname{fl}({\rm inner}(c))-M^{-1}c\right\rVert
       \lesssim\kappa u\left\lVert c\right\rVert\] and the explicit nonpathological-output assumption (6.4). The latter is an assumption, not a deterministic consequence of round-to-nearest arithmetic.\par
The SVD calculation in Appendix A.1 gives the componentwise backward-stability criterion \[|v_j^T(\widehat x-x)|
 \lesssim \sigma_j^{-1}(1+\left\lVert\widehat x\right\rVert)u
          +\sigma_j^{-2}\left\lVert b-A\widehat x\right\rVert u.\] The proof keeps the distinct \(\widehat R^{-1}e\) and residual-weighted \((A^TA)^{-1}e'\) terms through Lemma 6.7. The sketch-and-solve residual bound gives strong forward stability after one refinement; the second refinement and (6.4) give the stated backward criterion. This internal mechanism is consistent under the actual absolute preconditioner bounds. The exact-arithmetic sketch inequalities, the inverse perturbation estimates, and the SVD form of the Karlson–Waldén estimate were checked. External estimator and Lanczos theorems were not independently reproved.\par
\section*{Problems requiring correction}
Problem 1 (P01: Major—condition number alone does not suffice).  Theorem 6.10, p. 29, replaces sketch normalization by the condition that \(\operatorname{cond}(A\widehat R^{-1})\) is bounded, while retaining the absolute inner-error condition of Theorem 6.1. The resulting assertion is false.\par
For binary precision \(p\ge3\), put \(u=2^{-p}\) and choose \[A=b=\begin{pmatrix}1\\0\end{pmatrix},\qquad
 x=1,\quad \kappa=1,\quad
 \widehat R=u,\quad \widehat x_0=1+2u.\] Take the inner routine \(g(c)=uc\). Then \(M=u^{-2}\) and, for every scalar \(c\), \[|g(c)-M^{-1}c|=u(1-u)|c|\le\kappa u|c|.\] Thus hypothesis (D) holds. Hypothesis (A) holds as \(u\to0\) at fixed dimension. The triangular preconditioner is nonsingular and \(\operatorname{cond}(A\widehat R^{-1})=1\). The initial error satisfies Fact 6.9 with constant 2, since \(\left\lVert A(x-\widehat x_0)\right\rVert=2u\) and the exact residual is zero. All the data are representable binary numbers.\par
The first refinement is exact: \[r_0=-2u,\quad c_0=-2,\quad
 \delta y_0=-2u,\quad \delta x_0=-2,\quad
 \widehat x_1=-1+2u.\] Consequently \(|\widehat x_1-x|=2-2u\), which is not \(O(u)\). In the second refinement, \[r_1=2-2u,\quad c_1=2/u-2,\quad
 \delta y_1=2-2u,\quad \delta x_1=2/u-2.\] The exact sum for the updated iterate is \(2/u-3+2u\). Nearest-even \(p\)-bit rounding gives \[\widehat x_2=2/u-2.\] The output-norm lower bound (6.4) also holds: both iterates have norm bounded below by a fixed multiple of \(1+u\), and the second is large. Thus hypothesis (B) does not exclude the example.\par
The second iterate is not backward stable either. If perturbations satisfy \(\left\lVert\Delta A\right\rVert,\left\lVert\Delta b\right\rVert\le\epsilon<1\), the perturbed least-squares solution obeys \[|X|\le\frac{\left\lVert b+\Delta b\right\rVert}{\sigma_{\min}(A+\Delta A)}
      \le\frac{1+\epsilon}{1-\epsilon}.\] For \(X=\widehat x_2>1\), this requires \(\epsilon\ge(X-1)/(X+1)\), tending to one. Equality is attainable by scaling \(A\) down and \(b\) up along their common column. It therefore also fails the more restrictive \(A\)-only notion of backward stability used in the introduction.\par
Exact nearest-even checks at \(p=10,20,40,53\), together with an actual binary64 calculation, confirm every displayed operation and its ordinary relative-roundoff bound. These are independent reviewer calculations, not author-code executions. The analytical family, rather than those four sample precisions alone, establishes the counterexample.\par
The cause is amplification by \(\widehat R^{-1}\): bounded condition number does not fix the scale of \(A\widehat R^{-1}\). Add uniform two-sided absolute singular-value bounds, or state and prove a scale-aware inner-accuracy condition. The counterexample violates the absolute bounds supplied by sketch-QR and does not refute Theorem 6.1, the correctly normalized FOSSILS method, or the corresponding normalized SPIR mechanism.\par
Problem 2 (P02: Major—the recommended pseudocode is inconsistent).  Algorithm 4 on p. 19 is not an implementable statement of the recommended stable FOSSILS method as printed.\par
Parameter ratio. Lines 2–3 set \(d=12n\), \(\beta=d/n=12\) and \(\alpha=(1-\beta)^2=121\). The intended ratio is \(n/d\), consistent with \(\beta=\eta^2\) in (2.10). For a scalar normal matrix \(M=3/2\) and \(c=1\), with the prescribed two initial iterates equal to 1, the first printed update is \[1+121(1-\tfrac32)=-\tfrac{119}{2},
 \qquad M^{-1}c=\tfrac23.\] More generally the homogeneous error recurrence has characteristic root product \(\beta=12\), so at least one root has modulus greater than one. This nonzero initial-error example avoids the special case \(M=1\), in which \(y_0=y_1=c\) is already exact. Using \(\beta=1/12\), \(\alpha=121/144\), gives first update \(167/288\) and stable roots of modulus \(\sqrt{1/12}\) for the same scalar system.\par
Column scaling. Line 1 stores the column-norm vector in \(d\); line 2 overwrites it with the scalar sketch dimension. Line 37 then uses \(\operatorname{diag}(d)^{-1}\) to undo the lost scaling. Even for one column of norm 2, a correct normalized solution 2 must unscale to 1, whereas the overwritten value 12 gives \(1/6\). For several columns there is also a scalar/vector ambiguity.\par
Stopping test. Line 25 requires \(i=1\) and \(\operatorname{MOD}(i,5)=0\), although the outer index is \(i\in\{0,1\}\). The condition is never true. The intended periodic test must use the inner index \(j\).\par
Regularization. The \(+\mu^2z\) term in line 21 has an undefined parameter and is outside the projection. After truncation, \(P\) is \(n\times r\), while \(z=P\delta y\) is \(n\)-dimensional and \(c,\delta y\) are \(r\)-dimensional. The displayed addition is invalid when \(r<n\). A regularized normal operator must use the corresponding projected and scaled term, with the sign induced by \(c-M\delta y\).\par
These are errors in the supplied algorithmic prescription, including a noncontracting main iteration and incorrect output scaling. They are not attributed to the inspected code: `fossils.m` uses `size(c,1)/d`; `metasolver.m` keeps a distinct `scale` vector, uses the inner iteration in its periodic test, and applies the regularization inside the projected normal operator. The code is substantive counterevidence limiting this finding to the printed algorithm. Its numerical performance has not been independently certified.\par
Problem 3 (P03: Minor—the condition-number cap 40 does not follow).  The proof of Theorem 8.2, p. 35, invokes Fact 8.1 by asserting \(\operatorname{cond}(M)\le40\). The permitted distortion \(\eta\le0.9\) does not imply this value.\par
Take \(A=[I_n;0]\), \(\eta=7/8\), and \[SA=\operatorname{diag}(1/8,15/8,1,\ldots,1).\] This satisfies the stated embedding bounds, with exact QR and \(\kappa(A)=1\). The extreme eigenvalues of \(M\) are 64 and \(64/225\), so \(\operatorname{cond}(M)=225\). The dimension can be chosen to satisfy the additional logarithmic lower bound on \(n\), and the multiplication and factorization in this example are exact.\par
In general the exact upper bound is \[\operatorname{cond}(M)\le
       \left(\frac{1+\eta}{1-\eta}\right)^2\le361,\] with a small perturbation allowance for computed factors. The default \(\eta=1/2\) gives 9 and is consistent with cap 40. For the wider stated range, use a suitable fixed cap and adjust the fixed approximation interval in Appendix C. For any such fixed positive spectral interval, polynomial approximation of \(x^{-1}\) still gives the same qualitative logarithmic iteration order. This is a local numerical inference requiring repair in the application of the lemma; it is not a counterexample to stable SPIR output.\par
Problem 4 (P04: Minor—the Lanczos routine needs a zero-input branch).  Algorithm 6, p. 34, initializes \(\gamma=\left\lVert c\right\rVert\), \(q_1=c/\gamma\) without handling \(c=0\). Its later zero-\(\beta\) test cannot remedy the initial division.\par
The case is attainable in the outer method: an exact consistent sketch-and-solve start has \(\widehat x_0=x\) and \(c_0=0\). For example, take \(A=[I_n;0]\), \(b=Ae_1\) and an exact isometric sketch on the relevant range. Return the zero inner solution immediately when \(c=0\). This local boundary convention leaves the nonzero-input Lanczos analysis unchanged. No uninspected author Lanczos implementation is alleged to share the omission.\par
\section*{Additional checks and limits}
Theorem 8.2 specifically analyzes Lanczos linear solves. Exact-arithmetic equivalence with conjugate gradients or LSQR does not establish their floating-point equivalence. The manuscript itself presents the latter stability claims as empirical evidence and expectation; the review does not recast them as a proved theorem. The logarithmic-dimension assumption in Fact 8.1 is needed to absorb the polylogarithmic factors in Appendix C into the declared dimension-polynomial constants.\par
The proof of Lemma 7.3 contains a local initialization error: with \(d^{(j)}=M^{-1}c-y^{(j)}\) and \(y^{(0)}=y^{(1)}=c\), the initial error is \((M^{-1}-I)c\), not \(c\). The expressions involving \(j(1.01)^{j-1}\) on pp. 32–33 also omit the contracting factor \(\eta\), and a reference to assumption (C) in Lemma 6.7 should be to (D). These are recorded as local proof corrections; no false final Polyak bound is inferred from them. Uniformity of the \(1.01\eta\) inflation when \(\eta\) tends to zero with \(u\) was not resolved and is not promoted to a finding.\par
The review includes the definitions of forward and backward error, the componentwise characterization, normalized triangular-solve error structure, two-stage refinement, sketch-based error estimation, the stated Polyak and Lanczos scope, and the Prony example’s role as a least-squares application. It does not reproduce the timing comparisons, quantum data or benchmark plots. The four numbered problems comprise two Major and two Minor findings, with their version and dependency limits stated above. Human review and independent acceptance remain outstanding.\par
2 E. N. Epperly, M. Meier and Y. Nakatsukasa, Fast randomized least-squares solvers can be just as accurate and stable as classical direct solvers, arXiv:2406.03468v3, 2025. Task DOI: https://doi.org/10.1002/cpa.70013. Stable-Preconditioned-Solvers, retained source snapshot, commit `afc7c8c071e4347c22b729c52c8492aa5678f009`. Files `code/fossils.m` and `code/metasolver.m`; static reading only.\par
\end{document}
